\chapter{Page faults}
\label{CH:PGFAULTS}
%% 

The combination of page tables and page faults is a powerful tool for
kernels. It allows a kernel to transparent modify a user process's
page tables, which opens up a wide range of features.  It enables the
kernel to implement system calls more efficiently: for example,
copy-on-write fork allows the kernel to avoid copying the parents
memory. It can provide more convenience for applications programmers:
for example, lazy allocation allows an application to allocate memory
lazily without having to predict upfront exactly how memory it will
need.  xv6 doesn't make much use of this powerful combination; it has
only one example: lazy allocation.  This chapter describes this
example and provides examples of how production operating systems use
page-fault exceptions for other features.

Before proceeding, please read the functions {\tt sys\_sbrk()} in {\tt
  kernel/sysproc.c}, and {\tt vmfault} in {\tt kernel/vm.c}.

\section{Lazy allocation}
\label{sec:lazy}

A widely-used feature implemented using page faults is called
\indextext{lazy allocation}, which xv6 also implements. It has two
parts.  First, when an application asks for more memory by calling
\lstinline{sbrk}, the kernel notes the increase in size, but does not
allocate physical memory and does not create PTEs for the new range of
virtual addresses.  Second, on a page fault on one of those new
addresses, the kernel allocates a page of physical memory and maps it
into the page table.  The kernel can implement lazy allocation
transparently to applications: no modifications to applications are
necessary for them to benefit.

Lazy allocation is convenient for applications because they don't have
to accurately predict how much memory they will need.  For example, an
application may process input, but not know in advance how large the
input will be.  With lazy allocations an application can ask for
memory for the worst case, but not have to pay for this worst case:
the kernel doesn't have to do any work at all for pages that the
application never uses.

Furthermore, if the application is asking to grow the address space by
a lot, then \lstinline{sbrk} without lazy allocation is expensive: if
an application asks for a gigabyte of memory, the kernel has to
allocate and zero 262,144 4096-byte pages.  Lazy allocation allows
this cost to be spread over time.  On the other hand, lazy allocation
incurs the extra overhead of page faults, which involve a user/kernel
transition.  Operating systems can reduce this cost by allocating a
batch of consecutive pages per page fault instead of one page and by
specializing the kernel entry/exit code for such page-faults.

\section{Code}

The \lstinline{sbrk} system call shrinks or grows the memory of the
calling process. The call {\tt sbrk(n)} grows (or shrinks if {\tt n}
is negative) the process's memory size by n bytes, and then returns
the start of the newly allocated region (i.e., the old size).  The
system call is implemented by the function \lstinline{sys_sbrk}
\lineref{kernel/sysproc.c:/^sys_sbrk/}.  To grow the memory of a
process, {\tt sys\_sbrk} just increment the process's size
(\lstinline{myproc()->sz)) by {\tt n} and return the old size; it does
  not allocate memory.

When a process loads or stores to unallocated memory, the CPU will
raise \indextext{page-fault exception}, because the memory isn't
mapped by the page table, and thus the CPU cannot execute the
instruction.  \lstinline{usertrap} checks for this case
(\lineref{kernel/trap.c:/page fault/}) and calls \lstinline{vmfault}
(\lineref{kernel/vm.c:/^uvmfault/}) to handle the page fault.  It
checks that the fault address is a valid address in the faulting
process's memory and allocates physical memory with {\tt kalloc},
zeros the allocated memory, and adds PTEs to the user page table with
{\tt mappages}.  Xv6 sets the \lstinline{PTE_W}, \lstinline{PTE_R},
\lstinline{PTE_U}, and \lstinline{PTE_V} flags in the PTEs for the
allocated page. Then, \lstinline{usertrap} resumes the faulting
process at the instruction that caused the fault. Because the 
PTE now allows access, the re-executed load or store instruction
will execute without a fault.

If an application frees memory using {\tt sbrk}, \lstinline{sys_sbrk}
calls \lstinline{shrinkproc}, which calls \lstinline{uvmdealloc}.  The
real work is done by {\tt uvmunmap} \lineref{kernel/vm.c:/^uvmunmap/},
which uses {\tt walk} to find PTEs.  Since some page may never have
been used by the process and thus never have been allocated by {\tt
  vmfault}, {\tt uvmunmap} skips PTEs without the \lstinline{PTE_V}
flag.  If the PTE mapping is valid, {\tt uvmunmap} calls {\tt kfree}
to free the physical memory they refer to.

Note that Xv6 uses a process's page table not just to tell the
hardware how to map user virtual addresses, but also as the only
record of which physical memory pages are allocated to that
process. That is the reason why freeing user memory (in {\tt
  uvmunmap}) requires examination of the user page table.

\section{Copy-On-Write (COW) fork}

Many kernels use page faults to implement \indextext{copy-on-write
  (COW) fork}.  To explain copy-on-write fork, consider xv6's
\lstinline{fork}, described in Chapter~\ref{CH:MEM}.  \lstinline{fork}
causes the child's initial memory content to be the same as the
parent's at the time of the fork.  Xv6 implements fork with
\lstinline{uvmcopy} \lineref{kernel/vm.c:/^uvmcopy/}, which allocates
physical memory for the child and copies the parent's memory into it.
It would be more efficient if the child and parent could share the
parent's physical memory.  A straightforward implementation of this
would not work, however, since it would cause the parent and child to
disrupt each other's execution with their writes to the shared stack
and heap.

Parent and child can safely share physical memory by
appropriate use of page-table permissions and page faults.
The CPU raises an page-fault exception
when a virtual address is used that has no mapping
in the page table, or has a mapping whose \lstinline{PTE_V}
flag is clear, or a mapping whose permission bits
(\lstinline{PTE_R},
\lstinline{PTE_W},
\lstinline{PTE_X},
\lstinline{PTE_U})
forbid the operation being attempted.
RISC-V distinguishes three
kinds of page fault: load page faults (caused by load instructions),
store page faults (caused by store instructions),
and instruction
page faults (caused by fetches of instructions to be executed).  The
\lstinline{scause} register indicates the type of the
page fault and the \indexcode{stval} register contains the address
that couldn't be translated, which are passed to {\tt vmfault} in xv6.

The basic plan in COW fork is for the parent and child to initially
share all physical pages, but for each to map them read-only (with the
\lstinline{PTE_W} flag clear).
Parent and child can read from the shared physical memory.
If either writes a given page,
the RISC-V CPU raises a page-fault exception.
If xv6 were more sophisticated, {\tt vmfault}
can respond by allocating a
new page of physical memory and copying into it
the physical page that the faulted address maps to.
The kernel changes the relevant PTE in the faulting process's page
table to point to the copy and to allow writes as well as reads,
and then resumes the faulting
process at the instruction that caused the fault. Because the 
PTE now allows writes, the re-executed instruction
will execute without a fault.

Copy-on-write requires book-keeping
to help decide when physical pages can be freed, since each page can
be referenced by a varying number of page tables depending on the history of
forks, page faults, execs, and exits. This book-keeping allows
an important optimization: if a process incurs a store page
fault and the physical page is only referred to from that process's
page table, no copy is needed.

Copy-on-write makes \lstinline{fork} faster, since \lstinline{fork}
need not copy memory. Some of the memory will have to be copied
later, when written, but it's often the case that most of the
memory never has to be copied.
A common example is
\lstinline{fork} followed by \lstinline{exec}:
a few pages may be written after the \lstinline{fork},
but then the child's \lstinline{exec} releases
the bulk of the memory inherited from the parent.
Copy-on-write \lstinline{fork} eliminates the need to
ever copy this memory.
Furthermore, COW fork is transparent:
no modifications to applications are necessary for
them to benefit.

\section{Demand paging}

Yet another widely-used feature that exploits page faults is
\indextext{demand paging}.  In \lstinline{exec}, xv6 loads all 
of an application's text
and data into memory before starting
the application.  Since applications
can be large and reading from disk takes time, this startup cost can
be noticeable to users. To
decrease startup time, a modern kernel doesn't initially load
the executable file into memory, but just creates the user page table with
all PTEs marked invalid. The kernel starts the program running;
each time the program uses a page for the first time, a page
fault occurs, and in response
the kernel reads the content of the page from disk and
maps it into the user address space.  Like COW fork and lazy
allocation, the kernel can implement this feature transparently to
applications.

The programs running on a computer may need more memory than the
computer has RAM. To cope gracefully, the operating system may
implement \indextext{paging to disk}. The idea is to store only a
fraction of user pages in RAM, and to store the rest on disk in a
\indextext{paging area}. The kernel marks PTEs that correspond to
memory stored in the paging area (and thus not in RAM) as invalid. If
an application tries to use one of the pages that has been {\it paged
  out} to disk, the application will incur a page fault, and the page
must be {\it paged in}: the kernel trap handler will allocate a page
of physical RAM, read the page from disk into the RAM, and modify the
relevant PTE to point to the RAM.

What happens if a page needs to be paged in, but there is no free
physical RAM? In that case, the kernel must first free a physical page
by paging it out or {\it evicting} it to the paging area on disk, and
marking the PTEs referring to that physical page as invalid. Eviction
is expensive, so paging performs best if it's infrequent: if
applications use only a subset of their memory pages and the union of
the subsets fits in RAM. This property is often referred to as having
good locality of reference. As with many virtual memory techniques,
kernels usually implement paging to disk in a way that's transparent
to applications.

Computers often operate with little or no {\it free} physical memory,
regardless of how much RAM the hardware provides. For example, cloud
providers multiplex many customers on a single machine to use their
hardware cost-effectively. As another example, users run many
applications on smart phones in a small amount of physical memory. In
such settings allocating a page may require first evicting an existing
page. Thus, when free physical memory is scarce, allocation is
expensive.

Lazy allocation and demand paging are particularly advantageous when
free memory is scarce and programs actively use only a fraction of
their allocated memory. These techniques can also avoid the work
wasted when a page is allocated or loaded but either never used or
evicted before it can be used.

\section{Memory-mapped files}

Other features that combine paging and page-fault exceptions include
automatically extending stacks and \indextext{memory-mapped files},
which are files that a program mapped into its address space using
the \texttt{mmap} system call so that the program can read and write
them using load and store instructions.


\section{Real world}

Production operating systems implement copy-on-write fork, lazy
allocation, demand paging, paging to disk, memory-mapped files, etc.
Furthermore, production operating systems try to store
something useful in all areas of physical memory, typically caching
file content in memory that isn't used by processes.

Production operating systems provide applications with system
calls to manage their address spaces and implement their own
page-fault handling through the {\tt mmap}, {\tt munmap}, and {\tt
  sigaction} system calls, as well as providing calls to pin memory
into RAM (see {\tt mlock}) and to advise the kernel how an application
plans to use its memory (see {\tt madvise}).

%% 
\section{Exercises}
%% 

\begin{enumerate}

\item Write a user program that grows its address space by one byte by calling
\lstinline{sbrk(1)}.
Run the  program and investigate the page table for the program before the call
to
\lstinline{sbrk}
and after the call to
\lstinline{sbrk}.
How much space has the kernel allocated?  What does the
PTE
for the new memory contain?

\item Implement COW fork.

\item Implement {\tt mmap}.

\end{enumerate}
